\chapter{La gravità dell'algoritmo e la caverna digitale}

\begin{flushright}
	\textit{``I social media sono bias cognitivi con le gambe: si muovono, si riproducono, si evolvono.''}
\end{flushright}

\bigskip
\bigskip

\section{Il mito dell'algoritmo rotto}

Quante volte avete sentito qualcuno pronunciare questa frase, con un misto di frustrazione e rassegnazione: ``L'algoritmo \`{e} rotto, i social stanno distruggendo la nostra cultura''? \`{E} una sentenza che circola nei salotti, nelle redazioni dei giornali, nelle aule universitarie. \`{E} diventata quasi un mantra dell'intellettuale moderno, una scorciatoia comoda per spiegare il degrado della conversazione pubblica. Eppure, alla luce di tutto ci\`{o} che la scienza cognitiva ci ha insegnato sul cervello, questa spiegazione \`{e} profondamente sbagliata. O meglio: \`{e} una bugia confortante che ci assolve troppo in fretta.

L'algoritmo non \`{e} rotto. Funziona con una precisione chirurgica che farebbe invidia a qualsiasi orologio svizzero. La sua apparente tossicit\`{a} non nasce da un errore di programmazione nelle stanze della Silicon Valley, ma dallo sfruttamento sistematico, scientifico e spietato delle vulnerabilit\`{a} che l'evoluzione ha costruito dentro di noi nel corso di milioni di anni. Non \`{e} il software a essere difettoso. Siamo noi, nel senso pi\`{u} affascinante e inquietante del termine. Per capire davvero questo meccanismo dobbiamo fare un passo indietro e guardare la nostra specie dall'alto, come farebbe un naturalista che osserva un formicaio.

Dopo aver passato i capitoli precedenti a esplorare i nostri bias come reperti di un cervello antico, ci troviamo ora di fronte alla prima forza nella storia che quei bias non si limita a subirli: li studia, li misura, li alimenta. Per la prima volta l'animale spaventato che ci portiamo dentro ha trovato qualcuno disposto a sussurrargli all'orecchio tutto il giorno, sapendo esattamente cosa vuole sentirsi dire.

\section{La piramide di Wallace}

Lo scrittore americano David Foster Wallace, uno degli intelletti pi\`{u} acuti della sua generazione, aveva intuito qualcosa di fondamentale sulla natura umana prima ancora che i social esistessero nella forma che conosciamo. Aveva compreso una legge quasi fisica applicata alla sociologia: le persone vengono unite in massa dai loro istinti pi\`{u} bassi, e separate dai loro interessi pi\`{u} elevati. Immaginiamo allora la nostra specie come una gigantesca piramide. Non una piramide di potere o di ricchezza, ma di complessit\`{a} cognitiva.

Alla base si trova il nostro hardware biologico ancestrale, il regno di ci\`{o} che gli psicologi chiamano Sistema 1:\footnote{Kahneman, D. (2011). \textit{Pensieri lenti e veloci}. Mondadori, Milano.} quel processore automatico, rapidissimo e inconscio che gestisce la paura, il desiderio, la rabbia, la curiosit\`{a} verso il pericolo. A questo livello siamo tutti, magnificamente e disperatamente, identici. Non importa quanti anni abbiamo studiato, quante lingue parliamo, quanti libri abbiamo letto. Provate a immaginare un esperimento. Mettete davanti a uno schermo un premio Nobel per la fisica e un concorrente qualunque di un reality televisivo, e mostrate loro, senza preavviso, un incidente stradale, una rissa violenta, un contenuto fortemente sessualizzato. I loro cervelli reagiranno in modo quasi identico: una cascata di neurotrasmettitori attraverser\`{a} le stesse antiche strutture, l'amigdala si accender\`{a} come un allarme antincendio, il cortisolo e l'adrenalina faranno il loro ingresso puntuale.\footnote{LeDoux, J. (1996). \textit{The Emotional Brain}. Simon and Schuster, New York.} A questo livello della piramide, cultura, educazione e intelligenza sono sostanzialmente irrilevanti. La morbosit\`{a} e l'indignazione sono il pi\`{u} potente collante universale dell'umanit\`{a}.

Al vertice troviamo qualcosa di completamente diverso: l'intelletto raffinato, la neocorteccia al massimo del suo potenziale. \`{E} qui che vivono la poesia difficile, l'astrofisica, la filosofia, la musica da camera, la riflessione lenta e faticosa. Un territorio meraviglioso, ma con un difetto commerciale fatale: divide. La raffinatezza intellettuale ci frammenta in nicchie sempre pi\`{u} piccole e incomunicabili. Trovare dieci milioni di persone che si emozionino simultaneamente per un trattato di epistemologia \`{e}, statisticamente, quasi impossibile.

\begin{center}
	\captionof{figure}{La piramide cognitiva dell'umanit\`{a}: alla base gli istinti universali che ci uniscono, al vertice l'intelletto che ci distingue e ci separa. L'algoritmo punta sempre alla base, dove vive la massa.}
	\label{fig:piramide_wallace}
\end{center}

\section{La legge di gravità digitale}

\`{E} esattamente in questa sproporzione geometrica, in questo divario tra la base larghissima e il vertice strettissimo, che entra in gioco l'algoritmo con tutta la sua potenza. L'obiettivo primario di qualsiasi grande piattaforma non \`{e} mai stato l'educazione, n\'{e} il benessere degli utenti, n\'{e} il miglioramento del dibattito pubblico. L'obiettivo \`{e} uno solo, perseguito con una coerenza ammirevole nella sua brutalit\`{a}: catturare la nostra attenzione, trattenerla, monetizzarla.

Per creare un contenuto virale che raggiunga simultaneamente decine di milioni di persone, le piattaforme sono costrette, per pura necessit\`{a} matematica, a puntare alla base della piramide, al minimo comune denominatore cognitivo. Le reazioni di fronte a un saggio filosofico sono lente, disomogenee, imprevedibili. Quelle di fronte alla paura, allo scandalo, all'indignazione morale sono invece biologicamente garantite, veloci come un riflesso, universali come il battito cardiaco. L'algoritmo, in questo senso, non ci odia. Non ha un'agenda morale, non prova soddisfazione nel renderci pi\`{u} stupidi o pi\`{u} arrabbiati. Semplicemente, con la freddezza di un'equazione, salta a pi\`{e} pari la nostra razionalit\`{a} e va a parlare direttamente con l'animale spaventato che ci portiamo dentro da trecentomila anni. Esercita su di noi una costante forza di gravit\`{a} digitale, che tira sempre verso il basso, con pazienza infinita e senza mai stancarsi.\footnote{Harris, T. (2017). \textit{How Technology Hijacks People's Minds}. Center for Humane Technology, San Francisco.}

Pensateci come a una corrente sottomarina. In superficie l'oceano pu\`{o} sembrare calmo, ma sotto quella corrente vi spinge di continuo in una direzione precisa, e voi la sentite appena, e quasi non vi accorgete di essere stati trascinati lontano dalla riva. Il punto, ed \`{e} ci\`{o} che rende il fenomeno cos\`{i} sottile, \`{e} che ogni singola scrollata sembra una vostra libera scelta. Tecnicamente lo \`{e}. Ma la corrente ha gi\`{a} deciso in quale direzione \`{e} pi\`{u} facile nuotare.

\section{Lo spacciatore di dopamina}

Quel gesto di trascinare il dito verso il basso per aggiornare il feed non \`{e} affatto casuale. \`{E} identico al movimento di tirare la leva di una slot machine, e non per coincidenza: i progettisti delle piattaforme hanno studiato a fondo le sale da gioco di Las Vegas prima di disegnare i nostri telefoni.\footnote{Harris, T. (2016). How technology is hijacking your mind. \textit{Medium}.} Il segreto si chiama ricompensa variabile. A volte, scorrendo, trovate qualcosa di interessante; a volte no. \`{E} proprio questa incertezza, e non il premio in s\'{e}, ad attivare i circuiti della dipendenza con la stessa forza del gioco d'azzardo. Sean Parker, ex presidente di Facebook, lo ha ammesso senza giri di parole: sfruttavamo una vulnerabilit\`{a} della psicologia umana, lo sapevamo, e lo abbiamo fatto comunque.

Qui per\`{o} occorre sfatare un equivoco diffuso, perch\'{e} cambia tutto. La dopamina non \`{e} la molecola del piacere, come spesso si crede. \`{E} la molecola dell'anticipazione, del desiderio. Non vi fa sentire bene quando ottenete la ricompensa, vi fa bramare ossessivamente la prossima. Ecco perch\'{e} potete scorrere per ore sentendovi sempre pi\`{u} vuoti: state inseguendo un piacere che \`{e} sempre alla scrollata successiva, e che proprio per questo non arriva mai del tutto. \`{E} come avere in tasca una slot machine collegata direttamente ai centri del desiderio del cervello. E poi ci stupiamo che la gente non riesca a smettere di controllarla.

\section{Le camere dell'eco e l'ombra sulla parete}

Le camere dell'eco, quelle stanze immaginarie dove le voci rimbalzano all'infinito amplificandosi senza mai incontrare un contraddittorio, esistevano molto prima di internet. Ma erano limitate dalla geografia: il bar del paese, il circolo politico, la parrocchia. I social hanno trasformato quelle nicchie locali in bunker globali, dove milioni di persone possono confermarsi a vicenda le stesse idee.\footnote{Sunstein, C. R. (2017). \textit{\#Republic: Divided Democracy in the Age of Social Media}. Princeton University Press, Princeton.} L'algoritmo di raccomandazione amplifica l'effetto: invece di esporci per caso a opinioni diverse, come accadeva sfogliando un giornale, ci mostra sempre pi\`{u} contenuti simili a ci\`{o} che abbiamo gi\`{a} gradito. E le aziende lo sanno perfettamente. Documenti interni hanno rivelato la consapevolezza che certi algoritmi amplificano i contenuti pi\`{u} divisivi e rabbiosi, e la scelta di non intervenire, perch\'{e} l'indignazione genera coinvolgimento e il coinvolgimento genera profitti.\footnote{Haugen, F. (2021). Testimonianza alla Commissione Commercio del Senato degli Stati Uniti.} \`{E} come affidare la gestione del villaggio a qualcuno che guadagna ogni volta che scoppia una rissa.

A questo punto vale la pena fermarsi, perch\'{e} ci\`{o} che la tecnologia ha costruito non \`{e} affatto nuovo. \`{E} la versione elettrica della pi\`{u} antica immagine filosofica dell'illusione. Quasi venticinque secoli fa, nel settimo libro della \textit{Repubblica}, Platone immagin\`{o} degli uomini incatenati fin dalla nascita in fondo a una caverna, costretti a fissare una parete.\footnote{Platone. \textit{La Repubblica}, libro VII. Numerose edizioni.} Alle loro spalle un fuoco proietta ombre di oggetti che essi non possono vedere direttamente. Per quei prigionieri le ombre non sono una rappresentazione della realt\`{a}: sono l'unica realt\`{a} che conoscono. Hanno persino imparato a dare un nome alle ombre, a prevederne i movimenti, a competere su chi sappia riconoscerle pi\`{u} in fretta. Bastano poche modifiche per rendere quel mito perfettamente contemporaneo: la caverna \`{e} diventata lo schermo, il fuoco \`{e} la luce dei nostri dispositivi, e le ombre sono i contenuti che l'algoritmo proietta davanti a noi, selezionati uno per uno sulla base di ci\`{o} che ci tiene incatenati. Il prigioniero di Platone credeva di vedere il mondo; in realt\`{a} vedeva solo ci\`{o} che il dispositivo alle sue spalle gli lasciava vedere. La differenza, oggi, \`{e} che le nostre catene sono fatte di piacere anzich\'{e} di ferro, e che la parete si aggiorna in tempo reale per non annoiarci mai.

\section{Il bias della disponibilità e il mondo in fiamme}

Su queste pareti luminose, l'informazione pi\`{u} disponibile non \`{e} mai la pi\`{u} accurata, ma la pi\`{u} emotivamente coinvolgente. Le notizie false si diffondono molto pi\`{u} velocemente di quelle vere, perch\'{e} sono confezionate apposta per attivare i nostri istinti primitivi: paura, rabbia, indignazione.\footnote{Vosoughi, S., Roy, D., e Aral, S. (2018). The spread of true and false news online. \textit{Science}, 359(6380).} \`{E} come aver sostituito il telegiornale con un flusso personalizzato di tutto ci\`{o} che ci fa arrabbiare di pi\`{u}, convincendoci che quella sia una fotografia fedele del mondo.

C'\`{e} per\`{o} un meccanismo ancora pi\`{u} subdolo. Il cervello, lo abbiamo visto parlando di memoria, valuta la probabilit\`{a} di un evento in base alla facilit\`{a} con cui riesce a richiamarlo. \`{E} il bias della disponibilit\`{a}, quello che ci fa temere gli squali dopo aver visto un film di paura ambientato in mare. Ora moltiplicate questo effetto per un flusso infinito di drammi, scandali, disastri, crimini. Il risultato \`{e} che percepiamo un mondo infinitamente pi\`{u} pericoloso e corrotto di quanto sia in realt\`{a}. \`{E} il grande paradosso del nostro tempo: viviamo nell'epoca pi\`{u} sicura della storia umana, eppure ci sentiamo costantemente minacciati. Quando \`{e} stata l'ultima volta che avete letto una notizia banalmente positiva su un social? L'algoritmo ha imparato che la normalit\`{a} non fa cliccare, e che solo l'eccezionale, il drammatico, l'indignante riesce a bucare il rumore. \`{E} come vivere in un mondo dove esistono soltanto incendi, e mai i pompieri che li spengono.

\section{Il rabbit hole e il tribalismo planetario}

Alcune piattaforme video hanno poi perfezionato un'arte ulteriore, quella che in inglese si chiama \textit{rabbit hole}, la tana del coniglio. Iniziate guardando un innocuo video di gattini e, tre ore dopo, vi ritrovate convinti che la Terra sia piatta e che rettiliani governino il pianeta. Non \`{e} un caso. L'algoritmo di raccomandazione \`{e} ottimizzato per un solo obiettivo, tenervi incollati allo schermo, e ha scoperto che il modo migliore per riuscirci \`{e} condurvi per gradi verso contenuti sempre pi\`{u} estremi e polarizzanti. \`{E} il bias di conferma che incontra l'escalation algoritmica: ogni video rafforza la convinzione precedente e vi spinge un poco pi\`{u} in profondit\`{a} nella tana. Immaginate una libreria dove il libraio, ogni volta che comprate un libro di cucina, ve ne propone uno via via pi\`{u} radicale: prima ricette vegane, poi crudiste, infine volumi che sostengono che si possa vivere di sola aria. Tutto questo perch\'{e} guadagna a ogni minuto in cui restate dentro, indipendentemente da ci\`{o} che acquistate.

Questa discesa si innesta su una delle nostre eredit\`{a} pi\`{u} antiche, il tribalismo. Il cervello si \`{e} evoluto per gestire relazioni con circa centocinquanta persone, e oggi deve invece navigare reti di milioni.\footnote{Dunbar, R. I. M. (1992). Neocortex size as a constraint on group size in primates. \textit{Journal of Human Evolution}, 22(6).} Il risultato \`{e} che applichiamo logiche tribali primitive su scala planetaria: il ``noi contro loro'' non riguarda pi\`{u} il villaggio vicino, ma pu\`{o} significare miliardi di persone che ne detestano altri miliardi per ragioni sempre pi\`{u} astratte. Certe piattaforme sono strutturalmente progettate per il conflitto, con messaggi brevi che impediscono ogni sfumatura e con strumenti che premiano la condivisione indignata. Non a caso, le discussioni pi\`{u} feroci riguardano spesso argomenti di cui sappiamo pochissimo: \`{e} pi\`{u} facile odiare un'astrazione che una persona reale, e l'algoritmo ci tiene deliberatamente lontani dalla complessit\`{a}, perch\'{e} la complessit\`{a} non fa arrabbiare abbastanza.

\section{La realtà che si sgretola}

Prima dei social esisteva qualcosa che potremmo chiamare realt\`{a} condivisa. Tutti guardavano gli stessi pochi canali televisivi, leggevano gli stessi giornali, avevano accesso pi\`{u} o meno alle stesse informazioni. Si poteva dissentire sull'interpretazione, ma almeno i fatti di base erano comuni. Oggi ognuno vive nella propria bolla informativa personalizzata. Non interpretiamo soltanto la realt\`{a} in modo diverso: viviamo letteralmente in realt\`{a} diverse, perch\'{e} il vostro flusso \`{e} radicalmente distinto dal mio, calibrato sui vostri specifici bias.\footnote{Pariser, E. (2011). \textit{The Filter Bubble: What the Internet Is Hiding from You}. Penguin Press, New York.} \`{E} come se ciascuno abitasse un universo parallelo con le proprie leggi, la propria storia, i propri fatti. E poi ci stupiamo di non riuscire pi\`{u} a metterci d'accordo su nulla, e di incontrare persone intelligenti e colte convinte sinceramente di cose opposte sui fatti pi\`{u} elementari. Non \`{e} stupidit\`{a}: \`{e} il risultato di anni vissuti in ecosistemi informativi incompatibili.

Questa frammentazione non \`{e} soltanto un fastidio sociale, \`{e} una minaccia a qualcosa di pi\`{u} profondo. Negli anni Venti del secolo scorso, il giornalista Walter Lippmann aveva descritto un fenomeno che chiam\`{o} \textit{pseudo ambiente}: tra noi e il mondo reale, sosteneva, si frappone sempre una rappresentazione mentale, una mappa che ci costruiamo in testa e che non coincide mai del tutto con il territorio.\footnote{Lippmann, W. (1922). \textit{Public Opinion}. Harcourt, Brace and Company, New York.} Agiamo non sul mondo, ma su questa immagine del mondo. Lippmann scriveva pensando ai giornali e alla propaganda, ma la sua intuizione \`{e} diventata oggi una descrizione letterale: gli algoritmi hanno trasformato ogni utente in una camera stagna del proprio pseudo ambiente, costruito su misura, impermeabile a quello del vicino. Un secolo fa il rischio era che molti condividessero la stessa mappa distorta; oggi il rischio opposto \`{e} che non ne condividiamo pi\`{u} nessuna.

Ed \`{e} qui che la questione smette di essere psicologica e diventa politica. Il filosofo tedesco J\"urgen Habermas ha dedicato la vita a studiare ci\`{o} che chiamava la sfera pubblica: quello spazio ideale in cui cittadini liberi possono confrontarsi e deliberare insieme, lasciando che sia la forza del miglior argomento, e non quella del pi\`{u} forte, a orientare le decisioni comuni.\footnote{Habermas, J. (1981). \textit{Theorie des kommunikativen Handelns}. Suhrkamp, Francoforte.} Ma una deliberazione del genere ha una condizione di possibilit\`{a} elementare, che diamo per scontata solo finch\'{e} esiste: i partecipanti devono abitare la stessa realt\`{a}, condividere un minimo di fatti su cui ragionare. Si pu\`{o} discutere all'infinito su cosa sia giusto fare, ma non si pu\`{o} deliberare se non si concorda nemmeno su cosa stia accadendo. Frantumando la realt\`{a} condivisa in milioni di bolle incomunicabili, l'ecosistema digitale non avvelena soltanto il dibattito: ne erode le fondamenta stesse. Senza un mondo comune, la conversazione democratica non si svolge male, semplicemente non pu\`{o} pi\`{u} cominciare.

\section{Il pensiero come unica via d'uscita}

A questo punto sorge la domanda pi\`{u} importante di tutte: come si esce da questa macchina perfetta? La verit\`{a} scomoda \`{e} che non possiamo semplicemente disinstallare i nostri bias evolutivi, cablati nel cervello da centinaia di migliaia di anni. Ma c'\`{e} qualcosa di unico nella mente umana, una facolt\`{a} che nessun altro animale possiede al nostro grado: la capacit\`{a} di osservare i propri processi mentali mentre accadono. I filosofi e gli psicologi la chiamano metacognizione, letteralmente pensare al proprio pensiero. \`{E} insieme la nostra condanna e la nostra salvezza. Condanna, perch\'{e} ci rende dolorosamente consapevoli di quanto siamo inadeguati al mondo che abbiamo costruito. Salvezza, perch\'{e} una volta che vedi i fili del burattinaio, la magia perde una parte del suo potere.

Su questo punto vale la pena ascoltare una voce che ha attraversato il secolo pi\`{u} buio. La filosofa Hannah Arendt, riflettendo sulle catastrofi politiche del Novecento, giunse a una conclusione sorprendente: alla radice del male su larga scala non c'\`{e} quasi mai un'intelligenza diabolica, ma molto pi\`{u} spesso una semplice, diffusa incapacit\`{a} di pensare.\footnote{Arendt, H. (1978). \textit{The Life of the Mind}. Harcourt Brace Jovanovich, New York.} Non l'assenza di intelligenza, si badi, ma la rinuncia a fermarsi, a interrogarsi, soprattutto a guardare il mondo dal punto di vista di un altro. Per Arendt pensare significa anzitutto questo: la capacit\`{a} di mettersi mentalmente al posto di chi non la pensa come noi, di immaginare la prospettiva altrui. Ed \`{e} esattamente la facolt\`{a} che le bolle algoritmiche atrofizzano con pi\`{u} efficacia, perch\'{e} ci circondano solo di voci che ci somigliano e rendono l'altro invisibile, o peggio, una caricatura da disprezzare. Riconoscere questo pericolo non \`{e} un esercizio astratto: \`{e} comprendere che la prima vittima di una caverna su misura non \`{e} la verit\`{a} dei fatti, ma la nostra stessa capacit\`{a} di abitare un mondo insieme agli altri.

Cosa farne, in concreto? Come per il cibo, la diversit\`{a} \`{e} fondamentale. Seguire deliberatamente fonti che ci sfidano \`{e} scomodo come l'esercizio fisico, e altrettanto necessario: il cervello, come i muscoli, ha bisogno di resistenza per restare in forma. Non occorre per forza eliminare i social, anche se per qualcuno pu\`{o} essere la scelta migliore. Serve piuttosto coltivare una sorta di igiene digitale, creando zone franche senza dispositivi: la prima ora del mattino, i pasti, almeno un giorno alla settimana. \`{E} una forma di digiuno, ma applicata all'informazione, e concede al cervello quei periodi di silenzio di cui ha bisogno per consolidare i ricordi e rielaborare le esperienze. Periodicamente conviene fare una potatura delle proprie fonti, dolorosa ma necessaria come quella di un giardino. E soprattutto, prima di condividere qualunque contenuto che provochi una reazione emotiva violenta, vale la pena fermarsi un istante e chiedersi: sto reagendo con la parte pensante del cervello o con l'animale spaventato? \`{E} informazione, o \`{e} indignazione confezionata per generare un click?

C'\`{e} un'ultima via d'uscita, la pi\`{u} antica di tutte, e ce la indica ancora il mito della caverna. Per sottrarsi alla gravit\`{a} dell'algoritmo occorre un atto che va contro ogni istinto sociale: accettare la solitudine intellettuale come un prezzo non soltanto necessario, ma nobile. Spegnere il rumore di fondo, rifiutare l'indignazione preconfezionata, coltivare le proprie altezze cognitive pu\`{o} farci sentire, sul momento, terribilmente isolati, come se nuotassimo controcorrente mentre tutti gli altri si lasciano trasportare verso il basso con apparente piacere. Ma \`{e} in larga parte un'illusione ottica. Quella disconnessione dalla base della piramide non \`{e} la fine della vita sociale: \`{e} il primo passo del prigioniero che si volta verso la luce. Nel mito, chi spezza le catene e risale fuori dalla caverna soffre, accecato dal sole, e quando torna a raccontarlo agli altri viene deriso. Eppure ha visto. E se avremo la pazienza di sopportare questa temporanea solitudine, qualcosa accadr\`{a}: cominceremo a incontrare i nostri veri simili, non quelli che condividono i nostri riflessi primordiali, ma quelli che condividono le nostre domande, i nostri dubbi, le nostre passioni pi\`{u} autentiche. Smetteremo di essere scimmie intrappolate in un branco digitale, per tornare a essere, con fatica e con gioia, individui umani in compagnia. Quella compagnia, costruita lentamente, vale infinitamente di pi\`{u} di qualsiasi tendenza virale. Ed \`{e}, forse, l'unica forma di connessione che l'algoritmo non riuscir\`{a} mai a simulare.

\begin{center}
	\captionof{figure}{Una figura solitaria che nuota controcorrente in mezzo a una folla trascinata dalla corrente: l'immagine visiva del coraggio intellettuale nell'era dell'algoritmo.}
	\label{fig:individuo_folla}
\end{center}